\documentclass[11pt,a4paper]{article}
\usepackage[margin=2.2cm]{geometry}
\usepackage{amsmath,amssymb}
\usepackage{booktabs}
\usepackage{array}
\usepackage{xcolor}
\usepackage[colorlinks=true,linkcolor=blue,urlcolor=blue]{hyperref}
\usepackage{siunitx}

\newcommand{\tauh}{\tau_{h}}
\newcommand{\mtautau}{m_{\tau\tau}}
\newcommand{\mbb}{m_{bb}}
\newcommand{\mhh}{m_{HH}}
\newcommand{\ptmiss}{p_{T}^{\mathrm{miss}}}
\newcommand{\yes}{\textcolor{teal}{\ding{51}}}
\usepackage{pifont}
\newcommand{\ok}{\textcolor{teal}{$\checkmark$}}
\newcommand{\no}{\textcolor{red}{$\times$}}
\newcommand{\pt}{p_{T}}

\title{\vspace{-1.5cm}Delphes event selection vs.\ CMS HIG-25-008\\
\large what the ntuple converter applies, and what it cannot}
\author{}
\date{\today}

\begin{document}
\maketitle
\vspace{-0.8cm}

\begin{abstract}
\noindent
Comparison of the CMS Run-3 HH$\rightarrow b\bar b\tau\tau$ selection (HIG-25-008,
Table~1 and Sections~5--6) against the selection applied to the Delphes samples. The
purpose is to make every difference explicit, so that yields and sensitivities computed
on the Delphes samples are read with the right caveats rather than compared naively
against the published numbers.
\end{abstract}

\section{Scope}

The ``Delphes'' column throughout describes the HIG-25-008 resolved selection as
implemented in \texttt{scripts/delphes\_to\_sbi.py}, invoked with
\texttt{--cms-selection}. \textbf{It is not the default.} Without that flag the
converter applies a deliberately loose preselection --- one light lepton, one $\tauh$,
two jets, with no b-tag requirement, no charge requirement, no vetoes and no mass
windows --- which exists to mirror the inputs of the NSBI test rather than the CMS
analysis, and whose signal$/t\bar t$ acceptance ratio is correspondingly $\sim$2.5
rather than the hundreds a real $b\bar b\tau\tau$ selection gives. Numbers quoted from
Delphes should always state which of the two was used.

Both modes are semi-leptonic by default ($\tau_\mu\tauh$, $\tau_e\tauh$), matching the
scope of the NSBI test; \texttt{--channels mt,et,tt} adds $\tauh\tauh$, CMS's most
sensitive channel.

\section{Event content and channels}

\begin{table}[htbp]
\centering\small
\begin{tabular}{@{}lp{5.6cm}p{5.4cm}@{}}
\toprule
& \textbf{CMS HIG-25-008} & \textbf{Delphes} \\
\midrule
Channels & $\tau_\mu\tauh$, $\tau_e\tauh$, $\tauh\tauh$, boosted-$\tauh\tauh$
         & $\tau_\mu\tauh$, $\tau_e\tauh$ \ok; $\tauh\tauh$ opt-in \\
Channel assignment & $\mu \rightarrow \tau_\mu\tauh$, else $e \rightarrow \tau_e\tauh$,
                     else 2nd $\tauh \rightarrow \tauh\tauh$ (\S5)
                   & same priority \ok \\
Trigger & single-$\ell$, cross $\ell$-$\tau$, di-$\tau$, di-$\tau$+jet, HH-Parking,
          VBF, boosted (\S5) & \no\ none \\
Categories & boosted / VBF / resolved 1b1j / resolved 2b $\rightarrow$ 10 SRs (\S6)
           & \no\ none \\
Discriminant & ggHH and qqHH DNNs, 5-fold ensemble (\S6) & --- (NSBI replaces it) \\
\bottomrule
\end{tabular}
\end{table}

\section{Object selection}

\begin{table}[htbp]
\centering\small
\begin{tabular}{@{}lp{5.6cm}p{5.4cm}@{}}
\toprule
& \textbf{CMS (Table~1)} & \textbf{Delphes} \\
\midrule
Muon $\pt$ & $>22/26/6\,\si{\GeV}$ (cross $\mu$-$\tau$ / single $\mu$ / VBF+$\mu$)
           & $>22$ (cross) \ok \\
Muon $|\eta|$ & $<2.4$ & $\le 2.4$ \ok \\
Muon ID, isolation & Tight Muon ID, PF rel.\ iso $<0.15$ & \no\ none \\
\midrule
Electron $\pt$ & $>25/31/13(18)\,\si{\GeV}$ (cross $e$-$\tau$ / single $e$ / VBF+$e$)
               & $>25$ (cross) \ok \\
Electron $|\eta|$ & $<2.5$ & $\le 2.5$ \ok \\
Electron ID & Tight MVA Isolation+ID (80\% eff.) & \no\ none \\
\midrule
Lepton $|d_{xy}|$, $|d_z|$ & $<0.045$, $<0.2\,\si{\cm}$ & \no\ not available \\
\midrule
$\tauh$ $\pt$ & 32 ($\tau_\mu\tauh$), 35 ($\tau_e\tauh$),
                40/35/25 ($\tauh\tauh$, trigger-dependent)
              & 32 / 35 \ok; 40 ($\tauh\tauh$) \ok \\
$\tauh$ $|\eta|$, $|d_z|$ & $<2.5$, $<0.2\,\si{\cm}$ & $\le 2.5$ \ok, no $d_z$ \\
$\tauh$ identification & \textsc{DeepTau}: $D^\tau_{\mathrm{jet}}$ Medium,
                         $D^\tau_e$ VVLoose, $D^\tau_\mu$ VLoose/Tight
                       & \no\ one tag bit \\
$\tauh$ object & HPS; decay modes 1-prong$+0,1,2\pi^0$, 3-prong$+0,1\pi^0$
               & \no\ AK4 jet with a tag bit \\
\midrule
$H\rightarrow b\bar b$ jets & AK4, $\pt>20$, $|\eta|<2.5$,
                              $\Delta R(j,\ell)>0.5$ from both $\tau$ candidates (\S5)
                            & \ok\ exactly as CMS \\
b-jet pair assignment & HH-\textsc{bTag} DNN (\S5) & \no\ two highest b-tag \\
b-tagging & \textsc{UParT} (2024); categories split at the \textbf{10\% mistag} WP (\S6)
          & \texttt{--btag-min} at \textsc{UParT} Medium ($\sim$1\% mistag) \\
\bottomrule
\end{tabular}
\end{table}

\section{Pair selection and signal region}

\begin{table}[htbp]
\centering\small
\begin{tabular}{@{}lp{5.6cm}p{5.4cm}@{}}
\toprule
& \textbf{CMS} & \textbf{Delphes} \\
\midrule
Opposite charge & required (Table~1) & \no\ no --- see \S\ref{sec:cannot} \\
$\Delta R(\ell,\tauh)$ & $>0.5$ (resolved); $<0.8$ (boosted) & $>0.5$ \ok \\
Additional-lepton veto & required (Table~1, \S5) & yes \ok \\
Signal region & ellipse in $(\mtautau,\mbb)$, Eqs.~2--4 & yes \ok \\
$\mtautau$ used in the ellipse & DNN-regressed $\tau\tau$ mass (\S5)
  & FastMTT (closest available) \\
\bottomrule
\end{tabular}
\end{table}

\noindent The elliptical selections implemented, from Eqs.~2--4:
\begin{align*}
\frac{(\mtautau-116)^2}{61^2} + \frac{(\mbb-114)^2}{228^2} &< 1 \qquad (\tau_\mu\tauh) \\
\frac{(\mtautau-119)^2}{57^2} + \frac{(\mbb-109)^2}{232^2} &< 1 \qquad (\tau_e\tauh) \\
\frac{(\mtautau-117)^2}{54^2} + \frac{(\mbb-134)^2}{169^2} &< 1 \qquad (\tauh\tauh)
\end{align*}
CMS quotes $\sim$99\% signal efficiency for these, with background efficiencies of
62\%, 60\% and 54\% respectively.

\section{What Delphes cannot support, and why}
\label{sec:cannot}

\begin{description}\itemsep3pt
\item[Opposite charge.] The ntuple's \texttt{Jet} collection carried no charge field, so
  the $\tauh$ leg had no charge to compare and \texttt{--cms-selection} is charge-blind.
  This matters more than it sounds: same-sign pairs are CMS's QCD control region at
  $\sim$99\% purity (\S7), i.e.\ an essentially pure fake-$\tau$ population that CMS
  removes and we keep. \emph{Delphes does write \texttt{Jet.Charge}} --- it was simply not
  carried into the ntuple schema. Now added; it requires a re-ntuplisation (not a
  re-simulation) to take effect.
\item[Lepton ID, isolation, impact parameters.] Delphes leptons carry only the card's
  efficiency parameterisation. There is no ID discriminant, no isolation and no $d_{xy}$
  or $d_z$ to cut on.
\item[$\tauh$ object and identification.] A Delphes $\tauh$ \emph{is} an AK4 jet carrying
  a single tag bit, not an HPS-reconstructed visible $\tau$. Consequences: jet-level
  four-vector (underlying event inside $R=0.4$), no decay modes, and one bit in place of
  the three \textsc{DeepTau} discriminants. This shifts the whole $\tau\tau$ block ---
  $\mtautau$, $\Delta R_{\tau\tau}$, and the ellipse acceptance --- and no selection
  change addresses it.
\item[b-jet pair assignment.] CMS uses the HH-\textsc{bTag} DNN; we take the two
  highest-b-tag jets. This lands directly on $\mbb$ and $\Delta R_{bb}$.
\item[b-tag working point.] The Delphes bit is \textsc{UParT}-AK4 Medium (cut $0.1272$;
  card PATCH-6), i.e.\ $\sim$1\% mistag. CMS's resolved 1b1j/2b split uses the
  \textbf{10\% mistag} working point. \texttt{--btag-min 2} is therefore \emph{tighter}
  than CMS's 2b category boundary, not equivalent to it.
\item[Trigger and categorisation.] No trigger is applied, and the sample is flat rather
  than split into the ten category$\times$channel signal regions CMS fits.
\item[Boosted and VBF topologies.] AK8 jets and the VBF jet pair are not written to the
  ntuple, so those categories are absent entirely.
\end{description}

\section{Consequences worth stating alongside any number}

\begin{enumerate}\itemsep3pt
\item \textbf{The selection described here is not the converter's default.} It requires
  \texttt{--cms-selection}; the default preselection is much looser and must never be
  compared directly against published yields.
\item \textbf{$\mhh$ and $\mtautau$ are visible masses} in the converter output. The CMS
  reference ntuple takes $\mhh$ from a branch (\texttt{mass\_tautaubb}) whose definition
  should be confirmed --- if it is the corrected $\tau\tau$ system, the two are different
  observables and any $\mhh$ threshold behaves differently on each. FastMTT-corrected
  versions (\texttt{m\_hh\_fastmtt}, \texttt{m\_tautau\_fastmtt}) are carried alongside
  for exactly this reason.
\item \textbf{The stock (untuned) card underestimates the $j\rightarrow\tauh$ fake rate.}
  CMS quotes $10^{-2}$--$10^{-3}$ (\S7); in the untuned samples the fully-hadronic
  $t\bar t$ component contributes essentially nothing. The fake background is therefore
  \emph{too small}, and $S/B$ correspondingly optimistic. This is a detector-model
  deficiency, not a selection one, and is what the tuning $\tau$-mistag map addresses.
\item \textbf{$\Delta R(j,\ell)$ differs between the two modes} (0.4 vs 0.5). Small, but
  it lands on $\Delta R_{bb}$ and $\mbb$, both of which are training features.
\end{enumerate}

\vspace{0.4cm}
\noindent\footnotesize
Implementation: \texttt{scripts/delphes\_to\_sbi.py} ---
\texttt{cms\_select()} (CMS selection), \texttt{select()} (preselection),
\texttt{\_in\_ellipse()} (Eqs.~2--4), with the Table~1 thresholds and ellipse
coefficients pinned by tests in \texttt{tests/test\_delphes\_to\_sbi.py}.

\end{document}
